\section{Results}

\subsection{Matched constraints preserve complementary properties}
Score-only selection yields high IRES-LM scores but substantial parent drift
(Table~\ref{tab:ablation}). Energy, global-ensemble, and anchor terms each
improve their intended structural quantity. The combined
\textsc{StructIRES-Rank} preserves all three simultaneously: it reduces mean
absolute MFE deviation by 83.4\%, pairing-profile distance by 77.4\%, and
increases parent-anchor retention by 0.309 relative to score-only selection.
All arms operate on the same 512 candidates per parent and retain the same 50
candidates, so this pattern cannot be attributed to a larger search budget.

\begin{table*}[t]
\centering
\caption{Matched full-length IRES optimization ablation. Values are mean $\pm$
sample SD across 30 paired parent-by-pool units (ten parents and three frozen
pools). Every arm retains 50 variants from the same 512 candidates. The
direct-RNA proxy is held out from selection and is computational only.}
\label{tab:ablation}
\resizebox{\textwidth}{!}{%
\begin{tabular}{lccccc}
\toprule
Selection method & $|\Delta$MFE| & Pairing $L_1$ & State identity & Anchor retention & Direct-RNA proxy \\
\midrule
IRES-LM score-only & $4.021\pm1.254$ & $0.164\pm0.027$ & $0.798\pm0.048$ & $0.646\pm0.075$ & $0.419$ \\
+ energy rank & $0.647\pm0.196$ & $0.118\pm0.024$ & $0.857\pm0.042$ & $0.766\pm0.062$ & $0.427$ \\
+ ensemble rank & $2.089\pm0.666$ & $0.055\pm0.014$ & $0.932\pm0.054$ & $0.915\pm0.026$ & $0.431$ \\
+ anchor rank & $2.093\pm0.580$ & $0.062\pm0.015$ & $0.924\pm0.063$ & $0.920\pm0.028$ & $0.431$ \\
\textsc{StructIRES-Rank} & $0.666\pm0.226$ & $0.037\pm0.010$ & $0.954\pm0.049$ & $0.955\pm0.015$ & $0.435$ \\
\bottomrule
\end{tabular}}
\end{table*}

The independent held-out direct-RNA proxy increases from 0.419 for score-only
selection to 0.435 for the combined method. The paired difference is 0.0158
(95\% bootstrap interval 0.0079--0.0242); combined ranking also exceeds the
energy, ensemble, and anchor ablations by 0.0086, 0.0046, and 0.0045,
respectively. These results support complementary roles for the three
structural constraints, while remaining an in-silico enrichment result.

\begin{figure*}[t]
  \centering
  \includegraphics[width=0.98\textwidth]{figures/fig6_structires_ireslm_ablation.pdf}
  \caption{Matched \textsc{StructIRES-Rank} ablation. Each arm selects from
  the same candidate pool for every parent and random seed. Error bars are
  sample SD across 30 paired units. The direct-RNA proxy is held out from
  candidate selection and does not measure activity of selected variants.}
  \label{fig:ablation}
\end{figure*}

\subsection{Direct-RNA mutational knowledge acts as a transparent guardrail}
The public IAPV scan contains 208 A-substitution tiles covering 169 of 213
parent positions. For the IAPV-only matched ablation, adding its edit-risk term
decreases mean MPRA edit risk from $2.438\pm0.174$ to $1.782\pm0.087$
(paired difference $-0.656$, bootstrap interval $-0.933$ to $-0.480$). This
comes with a transparent MFE-preservation trade-off (0.641 to 0.887 kcal/mol),
so it is kept separate from the ten-parent headline result.

\begin{table}[t]
\centering
\caption{IAPV-only matched direct-RNA MPRA edit-risk guardrail. Both arms
select 50 variants from each of the same three 512-candidate pools. The risk
is induced by selected edits in observed parent mutational tiles; it is not a
new-candidate activity measurement.}
\label{tab:iapv}
\resizebox{\columnwidth}{!}{%
\begin{tabular}{lcccc}
\toprule
Method & Risk & $|\Delta$MFE| & Pairing $L_1$ & Anchors \\
\midrule
\textsc{StructIRES-Rank} & $2.438\pm0.174$ & $0.641\pm0.069$ & $0.032\pm0.006$ & $0.970\pm0.011$ \\
+ MPRA guardrail & $1.782\pm0.087$ & $0.887\pm0.131$ & $0.036\pm0.004$ & $0.963\pm0.008$ \\
\bottomrule
\end{tabular}}
\end{table}

\subsection{Assay shift motivates constrained rather than score-only design}
Released RNA-FM and UTR-LM checkpoint reruns recover their native benchmark
summaries, confirming that the published classifiers are meaningful baseline
proposal models. Yet on direct-RNA labels, RNA-FM checkpoint AUROCs range from
0.251 to 0.821 and the all-checkpoint ensemble has AUROC 0.600 and AUPR 0.058
at 5.1\% active prevalence. The released IRESfinder model is below random
ranking (AUROC 0.460) in the same transfer setting. Hence, the structural
terms are not presented as a substitute activity assay; they are conservative
constraints against over-optimizing an assay-specific score.
